% Part: sets-functions-relations
% Chapter: arithmetization
% Section: checking-details
%
\documentclass[../../../include/open-logic-section]{subfiles}


\begin{document}
	
\olfileid{sfr}{arith}{check}
\olsection{Geordende ringen en lichamen}
In dit hoofdstuk hebben we steeds gesteld dat bepaalde definities zich
``gedragen zoals het hoort''. In deze technische bijlage leggen we
precies uit wat we daarmee bedoelen en schetsen we hoe we kunnen
aantonen dat de definities zich inderdaad ``correct'' gedragen.
In \olref[int]{sec} definieerden we optelling en vermenigvuldiging op
$\Int$. We willen aantonen dat deze bewerkingen $\Int$, met die
definities, de structuur geven die we ervan ``verwachten''. Het gaat
hier in het bijzonder om de structuur van een commutatieve ring.
\begin{defn}
	Een \emph{commutatieve ring} is een verzameling $S$ met aangewezen elementen $0$ en $1$ en bewerkingen $+$ en $\times$, waarvoor de volgende acht formules gelden:
	\begin{align*}
		\emph{Associativiteit}&&a + (b+ c) & = (a + b) + c \\
		&& (a \times b) \times c & = a \times (b\times c)\\
		\emph{Commutativiteit}&&a + b &= b+ a  \\
		&&  a \times b&= b\times a\\ 
		\emph{Neutrale elementen}&&a + 0 &= a \\
		&& a \times 1 &= a\\
		\emph{Additieve inverse}&&(\exists b\in S)0&=a + b\\
		\emph{Distributiviteit}&&a \times (b+ c ) &= (a \times b) + (a \times c)
	\end{align*}
	Stilzwijgend zijn deze formules allemaal universeel gekwantificeerd, waarbij de variabelen alleen waarden in $S$ aannemen. Merk ook op dat de elementen $0$ en~$1$ hier niet de gelijknamige natuurlijke getallen hoeven te zijn.
\end{defn}
Om aan te tonen dat de gehele getallen een commutatieve ring vormen,
hoeven we dus alleen deze acht voorwaarden te verifiëren. Geen ervan
is moeilijk te bewijzen, maar alle controles samen kosten enig werk.
Zo bewijzen we de \emph{associativiteit} van de optelling als volgt:
\begin{proof} 
Neem $i, j, k \in \Int$. Er bestaan dan $a_1, b_1, a_2, b_2, a_3, b_3 \in
\Nat$ waarvoor $i = \equivrep{a_1, b_1}{}$ en $j =
\equivrep{a_2,b_2}{}$ en $k = \equivrep{a_3, b_3}{}$. (Om de
notatie leesbaar te houden, schrijven we ``$\equivrep{x, y}{}$'' in plaats van
``$\equivrep{x, y}{\Intequiv}$''; dit doen we in de hele
paragraaf.) Dan geldt:
\begin{align*}
	i + (j + k) &= \equivrep{a_1, b_1}{}+(\equivrep{a_2, b_2}{} + \equivrep{a_3, b_3}{}) \\
	&= \equivrep{a_1,  b_1}{} + \equivrep{a_2+a_3, b_2+b_3}{}\\
	&= \equivrep{a_1 + (a_2 + a_3), b_1 + (b_2 + b_3)}{}\\
	&= \equivrep{(a_1 + a_2) + a_3, (b_1 + b_2) + b_3}{}\\
	&= \equivrep{a_1 + a_2, b_1 + b_2}{} + \equivrep{a_3, b_3}{}\\
	&= (\equivrep{a_1, b_1}{} + \equivrep{a_2, b_2}{}) + \equivrep{a_3, b_3}{}\\
	&= (i+j) + k
\end{align*}
waarbij we zonder meer de eigenschappen van de optelling op $\Nat$
gebruiken.
\end{proof}
Op dezelfde manier bewijzen we de \emph{additieve inverse}:
\begin{proof}
Neem $i \in \Int$, zodat $i = \equivrep{a,b}{}$ voor zekere $a,b \in
\Nat$. Stel $j = \equivrep{b,a}{} \in \Int$. Uit de eigenschappen van
de natuurlijke getallen volgt $(a+b) + 0 = 0 + (b+a)$, zodat
$\tuple{a+b, b+a} \sim_\Int \tuple{0,0}$ krachtens de definitie, en dus
$\equivrep{a+b, b+a}{} = \equivrep{0, 0}{} = 0_\Int$. Bijgevolg is $i + j =
\equivrep{a,b}{}+\equivrep{b,a}{}=\equivrep{a+b, b+a}{}= \equivrep{0,
0}{} = 0_\Int$.
\end{proof}
En dit is een bewijs van de \emph{distributiviteit}:
\begin{proof}
Neem, zoals hierboven, $i = \equivrep{a_1, b_1}{}$ en $j =
\equivrep{a_2,b_2}{}$ en $k = \equivrep{a_3, b_3}{}$. Dan geldt:
\begin{align*}
	i \times (j + k) 
	&= \equivrep{a_1, b_1}{} \times (\equivrep{a_2,b_2}{} + \equivrep{a_3, b_3}{})\\
	&= \equivrep{a_1, b_1}{} \times \equivrep{a_2 + a_3,b_2+b_3}{}\\
	&= \equivrep{a_1  (a_2 + a_3) + b_1  (b_2+b_3), a_1  (b_2 + b_3) + b_1 (a_2 + a_3)}{}\\
	&= \equivrep{a_1 a_2 + a_1a_3 + b_1 b_2+b_1b_3, a_1 b_2 + a_1b_3 + a_2b_1 + a_3b_1}{}\\		
%		&= \equivrep{a_1 a_2 + b_1b_2 + a_1a_3 + b_1 b_2+b_1b_3, a_1 b_2 + a_1b_3 + a_2b_1 + a_3b_1}{}\\		
	&= \equivrep{a_1a_2 + b_1b_2, a_1b_2 + a_2b_1}{} + \equivrep{a_1a_3 + b_1b_3, a_1b_3 + a_3b_1}{}\\
	&= (\equivrep{a_1, b_1}{} \times \equivrep{a_2,b_2}{}) + (\equivrep{a_1, b_1}{} \times  \equivrep{a_3, b_3}{})\\
	&= (i \times j) + (i \times k)
\end{align*}
\end{proof}
Het bewijs van de overige vijf voorwaarden laten we als oefening.
Daarmee is aangetoond dat $\Int$ een commutatieve ring vormt, oftewel
dat optelling en vermenigvuldiging zich met de gegeven definities
gedragen zoals het hoort.
\begin{prob}
Bewijs dat $\Int$ een commutatieve ring is.
\end{prob}
Daarmee zijn we nog niet klaar. Naast optelling en vermenigvuldiging
op $\Int$ hebben we een ordeningsrelatie $\leq$ gedefinieerd; ook
daarvan moeten we controleren of die de gewenste eigenschappen
heeft. Preciezer: we moeten aantonen dat $\Int$ een \emph{geordende}
ring vormt.\footnote{Bedenk dat volgens
	\olref[sfr][rel][ord]{def:linearorder} een totale ordening
	een reflexieve, transitieve, antisymmetrische en samenhangende relatie is.
	Bij ordeningsrelaties wordt samenhangendheid soms
	\emph{trichotomie} genoemd, aangezien voor alle $a$ en $b$ geldt dat $a \leq b \lor a
	= b \lor a \geq b$.} 
\begin{defn}
Een \emph{geordende ring} is een commutatieve ring die bovendien is
voorzien van een totale ordening $\leq$, zodanig dat:
\begin{align*}
	a \leq b &\lif a + c \leq b + c\\
	(a \leq b \land 0 \leq c) &\lif a \times c \leq b \times c
\end{align*}
\end{defn}
\begin{prob}
Bewijs dat $\Int$ een geordende ring is. 
\end{prob}
Net als hierboven is het bewerkelijk maar routineus om aan te tonen
dat de geconstrueerde $\Int$ een geordende ring vormt. We laten dit
aan de lezer over.
Daarmee zijn de gehele getallen afgehandeld. Voor de rationale
getallen moeten we nu zeer vergelijkbare eigenschappen aantonen. We
moeten met name laten zien dat de rationale getallen een geordend
\emph{lichaam} vormen, uitgaande van onze definities van $+$,
$\times$ en $\leq$:
\begin{defn}\ollabel{orderedfield}
Een \emph{geordend lichaam} is een geordende ring waarvoor bovendien geldt:
\begin{align*}
	\emph{Multiplicatieve inverse}& & (\forall a \in S \setminus \{0\})(\exists b \in S) a\times b& = 1
\end{align*}
\end{defn}
Wanneer eenmaal is aangetoond dat $\Int$ een geordende ring vormt, is
het eenvoudig maar bewerkelijk om aan te tonen dat $\Rat$ een
geordend lichaam vormt.
\begin{prob}
Bewijs dat $\Rat$ een geordend lichaam is.
\end{prob}
Na de gehele en rationale getallen resteren alleen de reële getallen.
We moeten in het bijzonder aantonen dat $\Real$ een \emph{volledig}
geordend lichaam vormt, dat wil zeggen een geordend lichaam met de
volledigheidseigenschap. In \olref[cuts]{realcompleteness} is al
bewezen dat $\Real$ de volledigheidseigenschap heeft. We moeten echter
nog de (saaie) controles doorlopen die aantonen dat $\Real$ een
geordend lichaam is.
Voordat we aan \emph{die} bewerkelijke oefening beginnen, moeten we
nog enkele directere zaken controleren. Zo moeten we nagaan of
$\alpha + \beta$ volgens de definitie inderdaad een \emph{snede}
is, voor willekeurige sneden $\alpha$ en $\beta$. Dit bewijzen we als
volgt:
\begin{proof}	
Omdat $\alpha$ en $\beta$ beide sneden zijn, is $\alpha + \beta = \Setabs{p
+ q}{p \in \alpha \land q \in \beta}$ een niet-lege echte deelverzameling van
$\Rat$. Stel nu $x < p + q$ voor zekere $p \in \alpha$ en $q \in \beta$.
Dan is $x - p < q$, dus $x - p \in \beta$, en $x = p + (x - p) \in
\alpha + \beta$. Bijgevolg is $\alpha + \beta$ een beginstuk van $\Rat$.
Neem ten slotte een willekeurige $p + q \in \alpha + \beta$. Omdat
$\alpha$ en $\beta$ beide sneden zijn, bestaan er $p_1 \in \alpha$ en $q_1 \in \beta$
waarvoor $p < p_1$ en $q < q_1$; derhalve is $p + q < p_1 + q_1 \in \alpha +
\beta$; dus heeft $\alpha + \beta$ geen maximum.
\end{proof}
Op soortgelijke wijze kan worden aangetoond dat $\alpha - \beta$,
$\alpha \times \beta$ en $\alpha \div \beta$ sneden zijn (bij deling
met uitzondering van het geval waarin $\beta$ de nulsnede is). Ook
dit laten we aan de lezer over.
\begin{prob}
Bewijs dat $\Real$ een geordend lichaam is.
\end{prob}
Er rest nog een kleinigheid. In \olref[cuts]{sec} stellen we voor
$\sqrt{2} = \Setabs{p \in \Rat}{p < 0 \text{ of }p^2 < 2}$ te nemen.
We moeten echter wel aantonen dat deze verzameling een \emph{snede}
is. Dit bewijs luidt als volgt:
\begin{proof}
Dit is duidelijk een niet-leeg echt beginstuk van de rationale
getallen; het volstaat dus aan te tonen dat het geen maximum heeft.
Meer bepaald volstaat het aan te tonen dat voor elk positief
rationaal getal $p$ met $p^2 < 2$ en $q = \frac{2p+2}{p+2}$ zowel $p
< q$ als $q^2 < 2$ geldt. Voor $p < q$ geldt:
\begin{align*}
	p^2 &< 2\\
	p^2 + 2p &< 2 + 2p\\
	p(p + 2) &< 2 + 2p\\
	p &< \tfrac{2+2p}{p+2} = q
\end{align*}
Voor $q^2 < 2$ geldt:
\begin{align*}
	p^2 &< 2\\
	2p^2 + 4p + 2 &< p^2 + 4p+ 4\\
	4p^2 + 8p + 4 &< 2(p^2 + 4p + 4)\\
	(2p+2)^2 & <2(p+2)^2\\
	\tfrac{(2p+2)^2}{(p+2)^2} &< 2\\
	q^2 &< 2
\end{align*}
\end{proof}
\end{document}